% welcome-to-the-odyssey.tex
%
% Prompt 016: Welcome to the Reasoning Odyssey. Student-facing orientation
% deck explaining the shared academic fabric of A2 (the Reasoning Odyssey)
% -- NOT a weekly teaching deck, NOT world-specific lore. See
% sidecar/prompts/016_welcome_to_reasoning_odyssey.md, "Architectural
% firewall": this deck is layer 2 of 3 (class slides / THIS / world
% playgrounds). World names and one-line flavor below are drawn from the
% current source of truth (assignments/A2-coding-odyssey-project.md);
% no lore, characters, or hooks are invented here -- that is Prompt
% 017/018 territory.
%
% Built on the Prompt 015 CS2 slide system -- no second presentation
% style invented. See ../STYLE_GUIDE.md for the pattern vocabulary and
% the literal-frame convention required for patterns 5/6/7/8/13 (unused
% in this deck; no code or diagrams needed for an orientation talk).

\documentclass[aspectratio=169]{beamer}
\usepackage{../theme/cs2theme}
\usepackage{../theme/cs2commands}

\setcsweek{Orientation -- Welcome to the Odyssey}

\title{Welcome to the Reasoning Odyssey}
\subtitle{The shared spine underneath everything you'll build this semester}
\author{Computer Science II \textbullet\ Dr. Jeremy P. Evert}
\date{Fall 2026 \textbullet\ COMSC-1053}

\begin{document}

% 1. Title / opening
\begin{frame}[plain]
  \titlepage
\end{frame}

% 2. Today's destination
\begin{destinationframe}{Today's Destination}
  \begin{itemize}
    \item Understand what the Reasoning Odyssey actually is
    \item See what "one software world, all semester" means in practice
    \item Meet the World Bible and understand why it matters
    \item Understand how gates and checkpoints work
    \item Choose which world you'll build in
  \end{itemize}
  \instructorcue{This is the whole semester's shape in fifteen minutes. Slow down -- this is the only time you get to set expectations before the work starts.}
\end{destinationframe}

% 3. Why this matters
\begin{whymattersframe}{CS2 Is Not a Stack of Disconnected Assignments}
  Most programming courses hand you a new, unrelated exercise every week.
  Finish it, forget it, move on.

  \vspace{3mm}
  \textbf{This course does something different.} You are going to build
  \emph{one} software world, and every week's work either grows it,
  tests it, or forces you to defend a decision you made weeks ago.
\end{whymattersframe}

% 4. Core concept
\begin{conceptframe}{You Are Growing One World, Not Finishing Assignments}
  \begin{center}
    \large\color{cs2primary}
    \textbf{You are going to build one software world all semester.}
  \end{center}
  \vspace{2mm}
  It will start simple. Then we are going to stress it, break
  assumptions, add capabilities, and make you explain why it still
  deserves to exist.

  \vspace{3mm}
  \small That is not a metaphor -- it is literally how the weekly work
  is structured.
\end{conceptframe}

% 5. Design decision -- choices have consequences
\begin{designdecisionframe}{Your Choices Create Consequences}
  \begin{itemize}
    \item Early choices establish structure
    \item Later capabilities put pressure on earlier designs
    \item Assumptions may fail -- and that's expected, not a mistake
    \item You will refactor, collect evidence, and explain decisions
    \item An earlier design choice may matter again in Week 13
  \end{itemize}
  \vspace{2mm}
  \textbf{This continuity is a feature of the course, not an accident.}
\end{designdecisionframe}

% 6. World Bible entry
\begin{worldbibleframe}{The World Bible Remembers}
  \begin{worldbibleentry}{What the World Bible Actually Is}
    \wbfield{Purpose}{A short living project record: charter, current state, one line per week on what changed or broke, known debt, tests/evidence, and important design choices.}
    \wbfield{Not}{Decorative lore. Not busywork. Not a second homework track.}
    \wbfield{Why it matters}{When Week 13 exposes a consequence of a Week 5 decision, this is what you consult -- and what you cite in your defense.}
  \end{worldbibleentry}
\end{worldbibleframe}

% 7. Evidence to produce -- weekly gates
\begin{evidenceframe}{Weekly Gates Produce Evidence of Progress}
  Weeks 4--6, 8--9, and 11--14: a small, gradable evidence package, not a
  second problem-set track. Each gate asks for:
  \begin{enumerate}
    \item the world-fit implementation or design change
    \item a focused run, test, trace, or demonstration
    \item a short explanation of the decision and trade-off
    \item your World Bible entry, plus AI evidence when AI helped
  \end{enumerate}
  \vspace{2mm}
  \small 25 points each. Reasoning and explanation matter alongside working code.
\end{evidenceframe}

% 18. Odyssey checkpoint -- larger checkpoints
\begin{odysseygateframe}{Checkpoints Ask You to Defend a Larger Capability}
  Three larger checkpoints sit apart from the weekly gates -- bigger
  evidence packages, not just bigger point values:
  \begin{itemize}
    \item \textbf{Week 7} (40 pts) -- a baby-project contract and swap
    \item \textbf{Week 10} (50 pts) -- a GUI/model or data-story synthesis
    \item \textbf{Week 15} (60 pts) -- a professional workflow receipt and full trail debrief
  \end{itemize}
  \vspace{2mm}
  \small See \texttt{rubrics/odyssey\_gates/} for the full evidence shape --
  this slide is the pointer, not the rubric.
\end{odysseygateframe}

% 9. Shared destination
\begin{conceptframe}{Every World Shares the Same Destination}
  \begin{center}
    \large\color{cs2primary}\textbf{No world is the easy world.}\\
    \large\color{cs2primary}\textbf{No world is the advanced world.}
  \end{center}
  \vspace{3mm}
  Every world targets the same CS2 learning outcomes, the same gates,
  the same checkpoints, and the same grading model. You are choosing
  \emph{context and flavor} -- not a different standard.
\end{conceptframe}

% 10. Four doors, spoiler-light
\begin{frame}{Choose the World You'll Build In}
  \cskicker{cs2accent}{FOUR WORLDS, ONE DESTINATION}
  \begin{columns}[T]
    \begin{column}{0.24\textwidth}
      \begin{center}
        \textbf{\color{cs2primary}Frontier Settlement}\\[1mm]
        \small\color{cs2inkSoft} Build the services a new settlement needs to survive and grow.
      \end{center}
    \end{column}
    \begin{column}{0.24\textwidth}
      \begin{center}
        \textbf{\color{cs2primary}Investigation Bureau}\\[1mm]
        \small\color{cs2inkSoft} Build the casework and workflow tools an investigator actually relies on.
      \end{center}
    \end{column}
    \begin{column}{0.24\textwidth}
      \begin{center}
        \textbf{\color{cs2primary}Starship Log}\\[1mm]
        \small\color{cs2inkSoft} Build the systems that keep a ship's operations running -- and logged.
      \end{center}
    \end{column}
    \begin{column}{0.24\textwidth}
      \begin{center}
        \textbf{\color{cs2primary}Small Business}\\[1mm]
        \small\color{cs2inkSoft} Build the inventory and transaction tools a real small business needs.
      \end{center}
    \end{column}
  \end{columns}
  \vspace{4mm}
  \small\color{cs2inkSoft}\centering
  That's all you need to know today. The rest, you'll find out by building there.
\end{frame}

% 11. How to choose
\begin{pairdiscussframe}{How to Choose Without Overthinking It}
  Don't rank them by difficulty -- they're not ranked that way.

  \vspace{2mm}
  Ask yourself one question: \textbf{which of these four would I
  actually enjoy coming back to, week after week, for a whole semester?}

  \stopandwork{2 minutes. Talk to the person next to you about which door you're leaning toward -- and why.}
\end{pairdiscussframe}

% 12. What happens next
\begin{exitframe}{What Happens Next}
  \Large\color{cs2primary}
  \textbf{Pick your world. Week 4 is where the first gate opens.}

  \vspace{4mm}
  \normalsize\normalfont\color{cs2ink}
  Everything from here forward assumes you've chosen -- so choose
  something you'll want to keep building.
\end{exitframe}

\end{document}
